% textmakro.sty - Schnellreferenz
% Autor: OTS
% Datum: 28. Februar 2026

\documentclass[12pt,a4paper]{scrartcl}

% Pakete
\usepackage[ngerman]{babel}
\usepackage[utf8]{inputenc}
\usepackage[T1]{fontenc}
\usepackage{lmodern}
\usepackage{geometry}
\usepackage{xcolor}
\usepackage{tcolorbox}
\usepackage{booktabs}
\usepackage{longtable}
\usepackage{listings}
\usepackage{textmakro}

% Seitengeometrie
\geometry{
  left=2.5cm,
  right=2.5cm,
  top=2.5cm,
  bottom=2.5cm
}

% Listings-Konfiguration
\lstset{
  basicstyle=\ttfamily\small,
  breaklines=true,
  frame=single,
  backgroundcolor=\color{gray!10},
  keywordstyle=\color{blue},
  commentstyle=\color{green!60!black},
  stringstyle=\color{red},
  showstringspaces=false,
  numbers=left,
  numberstyle=\tiny\color{gray},
  language=[LaTeX]TeX
}

% Farbdefinitionen für Demonstration
\definecolor{verbcolor}{RGB}{50,205,50}
\definecolor{verbimperativcolor}{RGB}{255,255,0}
\definecolor{nomencolor}{RGB}{255,140,0}
\definecolor{konjunktioncolor}{RGB}{255,0,0}
\definecolor{ortcolor}{RGB}{255,0,255}
\definecolor{personcolor}{RGB}{100,100,205}
\definecolor{gnade}{RGB}{180,210,255}
\definecolor{bund}{RGB}{220,180,255}
\definecolor{gericht}{RGB}{255,180,180}

% Titel
\title{\textbf{textmakro.sty}\\[0.5em]\Large Schnellreferenz für Textmarkierungs-Makros}
\author{OTS}
\date{28. Februar 2026}

\begin{document}

\maketitle
\tableofcontents
\newpage

%%%%%%%%%%%%%%%%%%%%%%%%%%%%%%%%%%%%%%%%%%%%%%%%%%%%%%%%%
\section{Installation}
%%%%%%%%%%%%%%%%%%%%%%%%%%%%%%%%%%%%%%%%%%%%%%%%%%%%%%%%%

Um das \texttt{textmakro}-Paket zu verwenden, fügen Sie folgende Zeile in Ihre Präambel ein:

\begin{lstlisting}
\usepackage{textmakro}
\end{lstlisting}

Das Paket erfordert:
\begin{itemize}
  \item \texttt{colors.sty} -- Zentrale Farbdefinitionen
  \item \texttt{xparse} -- Für erweiterte Makro-Definitionen
\end{itemize}

%%%%%%%%%%%%%%%%%%%%%%%%%%%%%%%%%%%%%%%%%%%%%%%%%%%%%%%%%
\section{Verwendung}
%%%%%%%%%%%%%%%%%%%%%%%%%%%%%%%%%%%%%%%%%%%%%%%%%%%%%%%%%

Alle Makros in \texttt{textmakro.sty} unterstützen zwei Modi:

\begin{itemize}
  \item \textbf{Textfarbe:} \lstinline|\makro{Text}| -- Nur der Text wird eingefärbt
  \item \textbf{Box-Hintergrund:} \lstinline|\makro[b]{Text}| -- Text mit farbigem Hintergrund
\end{itemize}

\subsection{Syntax-Beispiel}

\begin{lstlisting}
\verb{ging}         % Text in Grün
\verb[b]{ging}      % Text mit grünem Hintergrund
\end{lstlisting}

%%%%%%%%%%%%%%%%%%%%%%%%%%%%%%%%%%%%%%%%%%%%%%%%%%%%%%%%%
\section{Makro-Übersicht}
%%%%%%%%%%%%%%%%%%%%%%%%%%%%%%%%%%%%%%%%%%%%%%%%%%%%%%%%%

%--------------------------------------------------------
\subsection{Verben}
%--------------------------------------------------------

\begin{longtable}{llp{8cm}}
\toprule
\textbf{Makro} & \textbf{Farbe} & \textbf{Verwendung} \\
\midrule
\endhead
%
\lstinline|\verb{Text}| & Grün (RGB 50,205,50) & Normale Verben \\
& & Beispiel: \lstinline|Abraham \verb{ging} nach Ur.| \\[0.5em]
%
\lstinline|\verbI{Text}| & Gelb (RGB 255,255,0) & Imperative Verben \\
& & Beispiel: \lstinline|\verbI{Gehe} und \verbI[b]{verkünde}!| \\[0.5em]
%
\lstinline|\verbP{Text}| & Grün + P & Passive Verben \\
& & Beispiel: \lstinline|Das Wort \verbP{wurde gesprochen}.| \\
\bottomrule
\end{longtable}

\subsubsection{Beispiele}

\begin{lstlisting}
Abraham \verb{ging} nach Ur.
\verbI{Gehe} und \verbI[b]{verkünde}!
Das Wort \verbP{wurde gesprochen}.
\end{lstlisting}

%--------------------------------------------------------
\subsection{Nomen}
%--------------------------------------------------------

\begin{longtable}{llp{8cm}}
\toprule
\textbf{Makro} & \textbf{Farbe} & \textbf{Verwendung} \\
\midrule
\endhead
%
\lstinline|\nomen{Text}| & Orange (RGB 255,140,0) & Nomen und Substantive \\
& & Beispiel: \lstinline|Das \nomen{Haus} des Herrn.| \\
\bottomrule
\end{longtable}

\subsubsection{Beispiel}

\begin{lstlisting}
Das \nomen{Haus} des \nomen[b]{Herrn}.
\end{lstlisting}

%--------------------------------------------------------
\subsection{Konjunktionen}
%--------------------------------------------------------

\begin{longtable}{llp{8cm}}
\toprule
\textbf{Makro} & \textbf{Farbe} & \textbf{Verwendung} \\
\midrule
\endhead
%
\lstinline|\konj{Text}| & Rot (RGB 255,0,0) & Bindewörter \\
& & Beispiel: \lstinline|Abraham \konj{und} Sara.| \\
\bottomrule
\end{longtable}

\subsubsection{Beispiel}

\begin{lstlisting}
Abraham \konj{und} Sara, \konj[b]{aber} auch Lot.
\end{lstlisting}

%--------------------------------------------------------
\subsection{Orte und Geographie}
%--------------------------------------------------------

\begin{longtable}{llp{8cm}}
\toprule
\textbf{Makro} & \textbf{Farbe} & \textbf{Verwendung} \\
\midrule
\endhead
%
\lstinline|\ort{Text}| & Magenta (RGB 255,0,255) & Geografische Orte \\
& & Beispiel: \lstinline|Von \ort{Jerusalem} nach \ort[b]{Bethlehem}.| \\
\bottomrule
\end{longtable}

\subsubsection{Beispiel}

\begin{lstlisting}
Von \ort{Jerusalem} nach \ort[b]{Bethlehem}.
\end{lstlisting}

%--------------------------------------------------------
\subsection{Personen -- Allgemein}
%--------------------------------------------------------

\begin{longtable}{llp{8cm}}
\toprule
\textbf{Makro} & \textbf{Farbe} & \textbf{Verwendung} \\
\midrule
\endhead
%
\lstinline|\person{Text}| & Blau (RGB 100,100,205) & Allgemeine Personen \\
& & Beispiel: \lstinline|\person{Abraham} war ein Glaubensmann.| \\[0.5em]
%
\lstinline|\mensch{Text}| & Blau (RGB 100,100,205) & Menschliche Wesen \\
& & Beispiel: \lstinline|\mensch{Adam} war der erste Mensch.| \\
\bottomrule
\end{longtable}

\subsubsection{Beispiel}

\begin{lstlisting}
\person{Abraham} war ein \mensch[b]{Mensch} des Glaubens.
\end{lstlisting}

%--------------------------------------------------------
\subsection{Theologische Personen -- Gottheit}
%--------------------------------------------------------

\begin{longtable}{llp{8cm}}
\toprule
\textbf{Makro} & \textbf{Farbe} & \textbf{Verwendung} \\
\midrule
\endhead
%
\lstinline|\gott{Text}| & Himmelsblau (RGB 180,210,255) & Gott \\
& & Beispiel: \lstinline|\gott{Gott} ist Liebe.| \\[0.5em]
%
\lstinline|\jesus{Text}| & Himmelsblau (RGB 180,210,255) & Jesus Christus \\
& & Beispiel: \lstinline|\jesus{Jesus} ist der Sohn Gottes.| \\[0.5em]
%
\lstinline|\geist{Text}| & Himmelsblau (RGB 180,210,255) & Heiliger Geist \\
& & Beispiel: \lstinline|Der \geist{Heilige Geist} wirkt.| \\
\bottomrule
\end{longtable}

\subsubsection{Beispiel}

\begin{lstlisting}
\gott{Gott} sandte \jesus[b]{Jesus}, und der \geist{Heilige Geist} wirkt.
\end{lstlisting}

%--------------------------------------------------------
\subsection{Theologische Personen -- Volk Gottes}
%--------------------------------------------------------

\begin{longtable}{llp{8cm}}
\toprule
\textbf{Makro} & \textbf{Farbe} & \textbf{Verwendung} \\
\midrule
\endhead
%
\lstinline|\israel{Text}| & Lila/Bund (RGB 220,180,255) & Israel \\
& & Beispiel: \lstinline|\gott{Gott} erwählte \israel[b]{Israel}.| \\
\bottomrule
\end{longtable}

\subsubsection{Beispiel}

\begin{lstlisting}
\gott{Gott} erwählte \israel[b]{Israel} als sein Volk.
\end{lstlisting}

%--------------------------------------------------------
\subsection{Theologische Personen -- Gegenspieler}
%--------------------------------------------------------

\begin{longtable}{llp{8cm}}
\toprule
\textbf{Makro} & \textbf{Farbe} & \textbf{Verwendung} \\
\midrule
\endhead
%
\lstinline|\teufel{Text}| & Rot/Gericht (RGB 255,180,180) & Teufel/Satan \\
& & Beispiel: \lstinline|\teufel{Der Teufel} versucht.| \\
\bottomrule
\end{longtable}

\subsubsection{Beispiel}

\begin{lstlisting}
\teufel{Der Teufel} versucht, \teufel[b]{Satan} widersetzt sich.
\end{lstlisting}

%--------------------------------------------------------
\subsection{Hilfsmakros}
%--------------------------------------------------------

\begin{longtable}{lp{10cm}}
\toprule
\textbf{Makro} & \textbf{Beschreibung} \\
\midrule
\endhead
%
\lstinline|\annot{Text}| & Annotation (grau, klein, kursiv) für Erklärungen \\
& Beispiel: \lstinline|\annot{Dies ist eine Erklärung.}| \\
\bottomrule
\end{longtable}

%%%%%%%%%%%%%%%%%%%%%%%%%%%%%%%%%%%%%%%%%%%%%%%%%%%%%%%%%
\section{Vollständiges Beispiel}
%%%%%%%%%%%%%%%%%%%%%%%%%%%%%%%%%%%%%%%%%%%%%%%%%%%%%%%%%

\begin{lstlisting}
\documentclass[12pt]{../../inc/mybib}
\author{OTS}

\setincpath{../../inc/}
\usepackage{bible_style}
\usepackage{textmakro}
\graphicspath{{../../assets/images/}}

\begin{document}

\section{Genesis 12,1-3}

\gott[b]{Der HERR} \verb{sprach} zu \person{Abram}: 

\verbI{Geh} aus deinem Land \konj{und} von deiner Verwandtschaft 
\konj{und} aus dem \nomen{Haus} deines Vaters in das \ort{Land}, 
das ich dir \verb{zeigen} werde.

\konj{Und} ich will dich zu einem gro{\ss}en \nomen{Volk} \verb{machen} 
\konj{und} will dich \verb{segnen} \konj{und} dir einen gro{\ss}en 
\nomen{Namen} \verb{machen}, \konj{und} du sollst ein \nomen{Segen} sein.

\annot{Dies ist der Anfang der Verhei{\ss}ung an Abraham.}

\end{document}
\end{lstlisting}

%%%%%%%%%%%%%%%%%%%%%%%%%%%%%%%%%%%%%%%%%%%%%%%%%%%%%%%%%
\section{Farbübersicht}
%%%%%%%%%%%%%%%%%%%%%%%%%%%%%%%%%%%%%%%%%%%%%%%%%%%%%%%%%

\begin{longtable}{lll}
\toprule
\textbf{Kategorie} & \textbf{Farbe} & \textbf{RGB} \\
\midrule
\endhead
%
Verben & Grün & rgb(50,205,50) \\
Imperativ & Gelb & rgb(255,255,0) \\
Nomen & Orange & rgb(255,140,0) \\
Konjunktionen & Rot & rgb(255,0,0) \\
Orte & Magenta & rgb(255,0,255) \\
Personen & Blau & rgb(100,100,205) \\
Gott/Jesus/Geist & Himmelsblau & rgb(180,210,255) \\
Israel & Lila & rgb(220,180,255) \\
Teufel & Rot (hell) & rgb(255,180,180) \\
\bottomrule
\end{longtable}

%%%%%%%%%%%%%%%%%%%%%%%%%%%%%%%%%%%%%%%%%%%%%%%%%%%%%%%%%
\section{Farben Deaktivieren (Boolean-Schalter)}
%%%%%%%%%%%%%%%%%%%%%%%%%%%%%%%%%%%%%%%%%%%%%%%%%%%%%%%%%

Alle Makros unterstützen Boolean-Schalter, um Farben selektiv zu deaktivieren. Dies ist nützlich für:

\begin{itemize}
  \item Schwarz-Weiß-Druck
  \item Vereinfachte Ansichten
  \item Selektive Hervorhebung einzelner Kategorien
\end{itemize}

%--------------------------------------------------------
\subsection{Verwendung}
%--------------------------------------------------------

Fügen Sie diese Zeilen \textbf{nach} \lstinline|\usepackage{textmakro}| in Ihrer Präambel ein:

\begin{lstlisting}
\ExplSyntaxOn

% Farben deaktivieren (false) oder aktivieren (true)
\bool_set_false:N \l_color_verb_bool      % Deaktiviert normale Verben
\bool_set_false:N \l_color_verbI_bool     % Deaktiviert Imperative
\bool_set_false:N \l_color_verbP_bool     % Deaktiviert Passiv-Verben
\bool_set_false:N \l_color_nomen_bool     % Deaktiviert Nomen
\bool_set_false:N \l_color_konj_bool      % Deaktiviert Konjunktionen
\bool_set_false:N \l_color_ort_bool       % Deaktiviert Orte
\bool_set_false:N \l_color_person_bool    % Deaktiviert Personen
\bool_set_false:N \l_color_gott_bool      % Deaktiviert Gott
\bool_set_false:N \l_color_jesus_bool     % Deaktiviert Jesus
\bool_set_false:N \l_color_geist_bool     % Deaktiviert Heiliger Geist
\bool_set_false:N \l_color_israel_bool    % Deaktiviert Israel
\bool_set_false:N \l_color_teufel_bool    % Deaktiviert Teufel

\ExplSyntaxOff
\end{lstlisting}

%--------------------------------------------------------
\subsection{Beispiel: Nur theologische Personen hervorheben}
%--------------------------------------------------------

\begin{lstlisting}
\documentclass[12pt]{../../inc/mybib}
\author{OTS}

\setincpath{../../inc/}
\usepackage{bible_style}
\usepackage{textmakro}

% Alle Farben au{\ss}er theologischen Personen deaktivieren
\ExplSyntaxOn
\bool_set_false:N \l_color_verb_bool
\bool_set_false:N \l_color_verbI_bool
\bool_set_false:N \l_color_verbP_bool
\bool_set_false:N \l_color_nomen_bool
\bool_set_false:N \l_color_konj_bool
\bool_set_false:N \l_color_ort_bool
\bool_set_false:N \l_color_person_bool
% Gott, Jesus, Geist, Israel, Teufel bleiben aktiviert (true)
\ExplSyntaxOff

\begin{document}
\gott{Gott} sprach zu Abraham...  % Nur Gott wird farbig sein
\end{document}
\end{lstlisting}

%--------------------------------------------------------
\subsection{Beispiel: Schwarz-Weiß-Druck}
%--------------------------------------------------------

\begin{lstlisting}
\ExplSyntaxOn
% Alle Farben deaktivieren
\bool_set_false:N \l_color_verb_bool
\bool_set_false:N \l_color_verbI_bool
\bool_set_false:N \l_color_verbP_bool
\bool_set_false:N \l_color_nomen_bool
\bool_set_false:N \l_color_konj_bool
\bool_set_false:N \l_color_ort_bool
\bool_set_false:N \l_color_person_bool
\bool_set_false:N \l_color_gott_bool
\bool_set_false:N \l_color_jesus_bool
\bool_set_false:N \l_color_geist_bool
\bool_set_false:N \l_color_israel_bool
\bool_set_false:N \l_color_teufel_bool
\ExplSyntaxOff
\end{lstlisting}

\begin{tcolorbox}[colback=yellow!10!white,colframe=orange!75!black,title=Hinweis]
Standardmäßig sind alle Farben aktiviert (\lstinline|\bool_set_true:N|). Sie müssen nur die Farben deaktivieren, die Sie nicht benötigen.
\end{tcolorbox}

%%%%%%%%%%%%%%%%%%%%%%%%%%%%%%%%%%%%%%%%%%%%%%%%%%%%%%%%%
\section{Tipps und Best Practices}
%%%%%%%%%%%%%%%%%%%%%%%%%%%%%%%%%%%%%%%%%%%%%%%%%%%%%%%%%

\begin{enumerate}
  \item \textbf{Konsistenz:} Verwenden Sie entweder nur Textfarben ODER nur Box-Farben in einem Dokument für ein einheitliches Erscheinungsbild.
  
  \item \textbf{Lesbarkeit:} Zu viele Box-Farben können überladen wirken -- setzen Sie sie sparsam ein.
  
  \item \textbf{Kombinationen:} Box-Farben eignen sich gut für wichtige Schlüsselwörter, während Textfarben für allgemeine Markierungen geeignet sind.
  
  \item \textbf{Annotationen:} Verwenden Sie \lstinline|\annot{}| für Erklärungen, die nicht Teil des Haupttextes sind.
  
  \item \textbf{Boolean-Schalter:} Nutzen Sie Boolean-Schalter, um Farben für Druck oder fokussierte Analysen zu deaktivieren.
  
  \item \textbf{Selektive Hervorhebung:} Deaktivieren Sie alle Farben au{\ss}er einer Kategorie, um spezifische grammatische oder theologische Aspekte zu analysieren.
\end{enumerate}

%%%%%%%%%%%%%%%%%%%%%%%%%%%%%%%%%%%%%%%%%%%%%%%%%%%%%%%%%
\section{Boolean-Schalter Referenz}
%%%%%%%%%%%%%%%%%%%%%%%%%%%%%%%%%%%%%%%%%%%%%%%%%%%%%%%%%

\begin{longtable}{lll}
\toprule
\textbf{Boolean-Variable} & \textbf{Makro(s)} & \textbf{Beschreibung} \\
\midrule
\endhead
%
\lstinline|\l_color_verb_bool| & \lstinline|\verb{}| & Normale Verben \\
\lstinline|\l_color_verbI_bool| & \lstinline|\verbI{}| & Imperative Verben \\
\lstinline|\l_color_verbP_bool| & \lstinline|\verbP{}| & Passive Verben \\
\lstinline|\l_color_nomen_bool| & \lstinline|\nomen{}| & Nomen/Substantive \\
\lstinline|\l_color_konj_bool| & \lstinline|\konj{}| & Konjunktionen \\
\lstinline|\l_color_ort_bool| & \lstinline|\ort{}| & Orte/Geographie \\
\lstinline|\l_color_person_bool| & \lstinline|\person{}, \mensch{}| & Personen \\
\lstinline|\l_color_gott_bool| & \lstinline|\gott{}| & Gott \\
\lstinline|\l_color_jesus_bool| & \lstinline|\jesus{}| & Jesus Christus \\
\lstinline|\l_color_geist_bool| & \lstinline|\geist{}| & Heiliger Geist \\
\lstinline|\l_color_israel_bool| & \lstinline|\israel{}| & Israel \\
\lstinline|\l_color_teufel_bool| & \lstinline|\teufel{}| & Teufel/Satan \\
\bottomrule
\end{longtable}

\end{document}
